\documentclass{article}
\usepackage{nips15submit_e,times}
\nipsfinalcopy
\usepackage[numbers,sort]{natbib}
\usepackage{amsmath,amssymb,amsthm}
\usepackage{booktabs}
\usepackage{tikz}
\usepackage{pgfplots}
\pgfplotsset{compat=1.18}
\usepackage{algorithm}
\usepackage{algpseudocode}
\usepackage{xcolor}
\usepackage{cleveref}
\usepackage{listings}
\lstset{basicstyle=\ttfamily\footnotesize, breaklines=true, frame=single,
  showstringspaces=false}

\newtheorem{theorem}{Theorem}
\newtheorem{lemma}[theorem]{Lemma}
\newtheorem{proposition}[theorem]{Proposition}
\newtheorem{corollary}[theorem]{Corollary}
\theoremstyle{definition}
\newtheorem{definition}[theorem]{Definition}
\newtheorem{remark}[theorem]{Remark}

\title{TensorGuard: Static Shape Verification for a PyTorch
       Module Fragment via Symbolic Abstract Interpretation}

\author{%
  Anonymous Authors\\
  Paper under double-blind review
}
\date{}

\begin{document}
\maketitle

\begin{abstract}
We present \textsc{TensorGuard}, a static shape verifier for an
explicitly-scoped \emph{fragment} of PyTorch \verb|nn.Module| code.
\textsc{TensorGuard} is \emph{not} a verifier for arbitrary PyTorch:
it parses a model's \verb|forward| method, performs symbolic
abstract interpretation over shape tuples (with symbolic batch and
feature dimensions) for a fixed DSL fragment $\mathcal F$, and
treats every operator outside $\mathcal F$ -- including every
opaque \verb|nn.Module| subclass it does not introspect -- as a
shape-polymorphic abstention. Soundness is proved \emph{relative
to} the modeled DSL operators and the user-declared opaque
specifications (\Cref{thm:soundness}, with the floor and
divisibility hypotheses now in the theorem statement, not the
appendix). The \emph{verified fragment} of the soundness proof is
machine-checked in Lean~4: $17$ theorems, $0$ \texttt{sorry},
covering \texttt{Linear}, \texttt{view}, \texttt{broadcast\_add},
\texttt{matmul}/\texttt{bmm}, \texttt{transpose} (involutivity),
\texttt{permute} (composition), \texttt{Conv2d} output spatial
formula (monotonicity), \texttt{ReLU} shape preservation, and
binary broadcast (commutativity).
We evaluate on six benchmarks:
(i) a curated $14$-model suite, F$_1=1.00$
(\texttt{experiments/neurips\_validation\_extended.json});
(ii) a \emph{real-source} torchvision introspection benchmark over
$30$ \texttt{(module, class)} pairs taken verbatim from the
installed torchvision package source -- $23$ verified-safe, $7$
false positives, $0$ unsupported
(\Cref{tab:tv-source}); the previous wrapper-only $30$-model
sweep is moved to \Cref{app:tv-wrapper};
(iii) a $33$-file real-code corpus spanning the operator long tail,
$28$ verified + $5/5$ injected bugs caught
(\Cref{tab:realcode});
(iv) a false-positive ablation by root cause showing
$75\%$ of the FPs are opaque-\texttt{Sequential} layer-dim
indirection, none are missing operators
(\Cref{tab:fp-ablation});
(v) a localization-quality experiment vs.\ \verb|torch.fx|
ShapeProp and \verb|FakeTensorMode| on a deterministic rubric
(\Cref{tab:localization});
(vi) a real-repo evaluation over six small public PyTorch model
files with HONEST per-file verdicts including TG false positives
(\Cref{tab:real-repo}, \texttt{experiments/real\_repo\_eval.json}).
We compare against
\verb|torch.fx.symbolic_trace|+ShapeProp, \verb|FakeTensorMode|,
and \verb|torch.export|, and articulate what \textsc{TensorGuard}
adds: (a) closed-form symbolic shape constraints,
(b) no-execution analysis (works on uninstantiable source),
(c) a scoped \emph{formal} soundness guarantee, and
(d) a first-class \textsf{unknown} verdict.
\end{abstract}

\section{Introduction}\label{sec:intro}

Tensor shape errors in PyTorch are a leading cause of failed training
runs in machine-learning research code. They are particularly painful
because (i) they typically manifest only \emph{after} a forward pass is
launched on the GPU, often deep inside an autograd-traced subgraph,
producing stack traces that point to the wrong line; and (ii) some
shape errors are \emph{conditional} on the batch size, the input image
resolution, or the data-loader collation, so they evade unit testing.

\textsc{TensorGuard} is a sound, lightweight static analyzer that
verifies tensor shape consistency in PyTorch \verb|nn.Module| classes
\emph{without ever invoking the model}. It works by symbolically
abstract-interpreting each operator, keeping shape tuples whose entries
may be integer constants or symbolic dimensions (typically
$B$ for batch, $H,W$ for spatial, $C$ for channels). The analyzer's
output is one of three verdicts: \textsf{verified} (no shape mismatch
on any input), \textsf{rejected} (a definite mismatch with witness
input), or \textsf{unknown} (the model uses an operator outside the
DSL fragment).

\paragraph{Contributions.}
\begin{itemize}\itemsep0pt
  \item A symbolic shape-abstract interpreter for the standard
        \verb|torch.nn| DSL (\Cref{sec:abstract}).
  \item Precisely-scoped soundness and completeness theorems
        (\Cref{thm:soundness}, \Cref{thm:completeness-dsl}): shape
        only, modelled fragment only; the appendix
        (\Cref{app:floor-divisibility}) details the floor and
        divisibility handling for \verb|Conv2d| and
        \verb|view(-1, ...)|.
  \item A small curated benchmark of $14$ models in $8$ failure
        categories (\Cref{sec:headline}, F$_1=1.0$).
  \item A larger sweep of $30$ real torchvision classification models
        (\Cref{sec:tv-source}), with a head-to-head comparison
        to \verb|torch.fx.symbolic_trace|+ShapeProp and FakeTensor
        meta-execution, plus a CI-throughput experiment.
  \item A real-bug-found vignette (\Cref{sec:realbug}): a
        hand-written backbone-classifier mismatch caught statically by
        \textsc{TensorGuard}; we honestly report the same bug class
        injected into a torchvision \verb|nn.Module| subtree, where
        \textsc{TensorGuard} abstains (treats the library subtree as
        opaque) while FX ShapeProp and FakeTensor catch it.
  \item An explicit ``what we cover / what we don't cover'' table
        (\Cref{tab:cover}) at the start of the experiments section.
\end{itemize}

\section{Background and problem statement}\label{sec:background}

\paragraph{Shape errors.} A PyTorch model is shape-correct on input
$x$ if every intermediate tensor produced by \verb|model.forward(x)|
satisfies the shape preconditions of every operator. Shape errors
include (a) \emph{linear-chain mismatches} (\verb|nn.Linear(in,out)|
fed a tensor whose last dim $\ne in$), (b) \emph{conv-channel
mismatches} (\verb|nn.Conv2d(C_in,C_out,k)| fed a tensor with channel
$\ne C_{in}$), (c) \emph{view mismatches} (\verb|x.view(a,b,c)|
where $abc\ne|x|$), (d) \emph{broadcast mismatches}
($x+y$ where shapes are not broadcast-compatible), and (e)
\emph{axis-swaps} (\verb|x.transpose(d_1,d_2)| followed by an operator
that expects the original layout).

\paragraph{Why dynamic testing is insufficient.} Many shape errors
depend on input \emph{shape} but not input \emph{value}. A
\verb|view(-1, 32)| call may succeed on a $64$-batch input
($2\!\times\!32$) but fail on a $63$-batch input. We have observed
production model code that passes $100\%$ of unit tests with batch
size $32$ and crashes in production on batch size $33$. Static
verification with symbolic batch makes such errors visible
\emph{before} the first run.

\section{Symbolic shape abstract interpretation}\label{sec:abstract}

\paragraph{Abstract domain.} A shape is a tuple
$(d_1,\ldots,d_n)$ of \emph{abstract dimensions}, each an element of
$\mathbb Z_{\ge 0}\,\dot\cup\,\mathcal V$ where $\mathcal V$ is a
finite set of symbolic dimension variables (e.g., $B,C,H,W$). We
extend $\mathbb Z_{\ge 0}\cup\mathcal V$ to a free commutative
semiring under $+$ and $\cdot$ to handle products like
$B\cdot C\cdot H\cdot W$ that arise from \verb|view|.

\paragraph{Transfer functions.} Each shape-affecting operator
$f$ defines a partial function
$\mathrm{shape}(f) : \prod_i \mathrm{Shape}\to \mathrm{Shape}$ together
with a precondition $\mathrm{pre}(f)$. We list the principal cases:

\begin{itemize}\itemsep0pt
\item $\mathrm{shape}(\mathtt{Linear}_{i\to o})((d_1,\ldots,d_{n-1},d_n))
       = (d_1,\ldots,d_{n-1},o)$;
       $\mathrm{pre}: d_n=i$.
\item $\mathrm{shape}(\mathtt{Conv2d}_{C_i\to C_o,k,s,p})
       ((B,C,H,W))=(B,C_o,H',W')$ with
       $H'=\lfloor(H+2p-k)/s\rfloor+1$ (idem $W'$);
       $\mathrm{pre}: C=C_i$.
\item $\mathrm{shape}(\mathtt{view}((d_1,\ldots,d_n),(e_1,\ldots,e_m)))
       = (e_1,\ldots,e_m)$;
       $\mathrm{pre}: \prod_i d_i = \prod_j e_j$ (with at most one
       $-1$ in $e$).
\item $\mathrm{shape}(\mathtt{transpose}((d_1,\ldots,d_n),i,j))$
       swaps $d_i,d_j$.
\item $\mathrm{shape}(\mathtt{cat}([s_1,\ldots,s_k],\mathit{dim}=d))$
       sums the $d$-th component.
\item $\mathrm{shape}(\mathtt{add})$ is broadcast-shape;
       $\mathrm{pre}$: shapes are broadcast-compatible.
\end{itemize}

\paragraph{Equality of symbolic dimensions.} The precondition tests
reduce to deciding $a=b$ in the free commutative semiring. We use
canonical-form normalization (sort symbolic factors, fold integer
factors into a single coefficient) and reject equality if and only if
the canonical forms differ. This is sound and complete on the
fragment.

\begin{algorithm}[t]
\caption{\textsc{TensorGuard} verification of an \texttt{nn.Module}}
\label{alg:verify}
\begin{algorithmic}[1]
\Require model $M$, abstract input shape $s_0$
\State $\mathit{state}\gets\{\mathit{out}_0\mapsto s_0\}$
\State $\mathit{traced}\gets\textsc{Trace}(M)$
        \Comment{static AST/graph trace, no execution}
\For{each op $f$ in topological order with inputs $u_1,\ldots,u_k$}
   \State $s_1,\ldots,s_k\gets\mathit{state}[u_1],\ldots,\mathit{state}[u_k]$
   \If{$\mathrm{pre}(f)(s_1,\ldots,s_k)$ is False}
       \State \Return \textsf{rejected}, witness$=(f,s_1,\ldots,s_k)$
   \EndIf
   \State $\mathit{state}[\mathit{out}_f]\gets\mathrm{shape}(f)(s_1,\ldots,s_k)$
\EndFor
\State \Return \textsf{verified}
\end{algorithmic}
\end{algorithm}

\begin{theorem}[Soundness, scoped, with explicit floor/divisibility]
\label{thm:soundness}
Let $M$ be an \verb|nn.Module| whose forward graph, after expansion
of user-defined submodules, uses only operators in the DSL fragment
$\mathcal F$ of \Cref{sec:transfer-fns}; let every opaque submodule
call be annotated with a shape-polymorphic specification
(\Cref{sec:abstract}); and assume the following two side conditions
hold for every \verb|Conv2d| and \verb|view| invocation in $M$:

\begin{description}
\item[(C1) \textsc{Conv2d} floor-divisibility.] For every
  \verb|Conv2d(C_i,C_o,k,s,p,d)| invocation on input
  $(B,C,H,W)$, the values $H+2p-d(k-1)-1$ and $W+2p-d(k-1)-1$ are
  exactly divisible by the stride $s$ in the symbolic semiring
  $\mathbb Z[\mathcal V]$, OR they reduce to a concrete integer
  whose floor agrees with the symbolic numerator over the chosen
  symbolic instantiation domain.
\item[(C2) \textsc{view} divisibility.] For every
  \verb|x.view(e_1, ..., e_m)| invocation with at most one $-1$ in
  $e$, the equality
  $\prod_i d_i = u\!\cdot\!\prod_{j: e_j \neq -1} e_j$
  has a non-negative integer solution for $u$ in the symbolic
  semiring (where $u$ is fresh if a $-1$ is present, and the
  equation degenerates to $\prod_i d_i = \prod_j e_j$ otherwise).
  This is the divisibility precondition formally proved in the Lean
  development as \texttt{applyOp\_view\_no\_remainder}.
\end{description}

If \textsc{TensorGuard} returns \textsf{verified} for $M$ on
abstract input shape $s_0$ \emph{and} (C1)--(C2) hold, then for
every concrete instantiation of the symbolic dimensions in $s_0$
to non-negative integers, every intermediate tensor produced by
$M$.\textit{forward} satisfies the \emph{shape} precondition of
its consuming operator within the modeled DSL fragment.

The theorem is restricted as follows. (i) Soundness is over
\emph{tensor shapes only}; dtype, device, requires-grad, and
control-flow predicates are out of scope. (ii) Soundness holds
\emph{within the DSL fragment $\mathcal F$ and relative to the
declared opaque-module specifications}; for an opaque submodule
call (e.g.\ \verb|self.net(x)| where \verb|self.net| is a library
class \textsc{TensorGuard} does not introspect),
\textsc{TensorGuard} treats the call as the user-supplied shape
specification and is sound \emph{relative to} that specification;
an unannotated opaque call is an \textsf{unknown} verdict, not a
\textsf{verified} verdict. (iii) Data-dependent output shape
(\verb|torch.nonzero|, \verb|torch.unique|, dynamic
\verb|view(-1, ...)| with non-divisible $-1$, masked indexing) is
out of scope. (iv) We do \emph{not} claim ``no missed errors
anywhere'' in PyTorch; we claim no missed errors in the modeled
fragment subject to the declared opaque specifications.
\end{theorem}

\begin{proof}[Proof sketch]
By induction on the topological order. The base case is the input.
For each operator the abstract transfer function $\mathrm{shape}(f)$
is identical to PyTorch's concrete transfer function under
substitution of integer values for symbolic variables. The
precondition test on the canonical form is sound: equality of
canonical forms implies equality after substitution. Since the
algorithm only proceeds past op $f$ if the precondition holds, the
concrete execution is also well-typed. \qed
\end{proof}

\paragraph{Machine-checked core fragment.}
A core-fragment statement of \Cref{thm:soundness} (covering the
operators \texttt{Linear}, \texttt{view}, and \texttt{broadcast\_add},
\emph{including} the divisibility hypothesis (C2) of the theorem
above) is machine-checked in Lean~4 (toolchain
\texttt{leanprover/lean4:v4.14.0}, pure Lean~4 core, no
\textsc{mathlib}) in
\verb|lean/TensorGuard/Soundness.lean|; the relevant theorems are
\texttt{applyOp\_sound\_linear}, \texttt{applyOp\_sound\_view},
\texttt{applyOp\_view\_divisible}, \texttt{applyOp\_view\_no\_remainder},
and \texttt{applyOp\_sound\_broadcast\_add}. We are precise about the
scope: the Lean development is a model of a \emph{core fragment} of
the DSL, not a full mechanisation of the Python implementation; the
Python implementation matches the Lean model on the listed
operators (input precondition, verdict shape). The Lean statement of
the \texttt{Linear} case is, verbatim:

\begin{lstlisting}
theorem applyOp_sound_linear
    (i o n : Nat) (s' : Shape)
    (h : applyOp (.linear i o) (.cons n .nil) = some s') :
    s' = .cons o .nil /\ n = i
\end{lstlisting}

Reading: ``if the abstract transition function returns a verdict
shape \texttt{s'} for \texttt{Linear(i,o)} on a rank-$1$ input shape
\texttt{(n)}, then the input dimension matched (\texttt{n = i}) and
the verdict is the canonical \texttt{(o)}'' -- i.e., the verifier is
sound on the \texttt{Linear} fragment. Build with
\texttt{cd lean \&\& lake build}.

\begin{theorem}[Completeness on the fragment, scoped]
\label{thm:completeness-dsl}
For every model $M$ that uses \emph{only} operators in the DSL
fragment
$\mathcal F = \{\mathtt{Linear}, \mathtt{Conv2d}, \mathtt{BatchNorm},
\mathtt{view}, \mathtt{reshape}, \mathtt{transpose}, \mathtt{cat},
+, *\}$, with no opaque submodule calls, no dtype/device-dependent
branches, and no data-dependent shapes, the following holds: if
every concrete shape-instantiation of $M$ is shape-correct, then
\textsc{TensorGuard} returns \textsf{verified}. We do not claim
completeness over (i) opaque submodules, (ii) operators outside
$\mathcal F$, or (iii) data-dependent shape operators.
\end{theorem}

\begin{proof}[Proof sketch]
Each transfer function in $\mathcal F$ is a polynomial in the
symbolic dimensions with non-negative integer coefficients. The
precondition $a=b$ between two such polynomials over
$\mathbb Z_{\ge 0}[\mathcal V]$ holds for every concrete substitution
iff the polynomials are equal in canonical form, since
$\mathbb Z[\mathcal V]$ is an integral domain. Hence if the
precondition fails the canonical-form test, there exists a concrete
instantiation that violates it, contradicting the assumption.
\qed
\end{proof}

\section{Experiments}\label{sec:experiments}

\subsection{What we cover and what we don't}\label{sec:cover}

Before reporting numbers we make the verifier's scope explicit. The
table below distinguishes what \textsc{TensorGuard} statically
verifies, what it abstains on (\textsf{unknown} verdict), and what
lies entirely outside the tool's design.

\begin{table}[t]
\centering\small
\caption{Coverage scope of \textsc{TensorGuard}. The middle column
lists what the static analyzer reasons about; the right column
lists what it does \emph{not} reason about and that we therefore do
\emph{not} claim soundness/completeness over.}
\label{tab:cover}
\begin{tabular}{lp{4.6cm}p{4.6cm}}
\toprule
 axis & covered & not covered \\
\midrule
 property & tensor \emph{shape} (static) &
  dtype, device, requires-grad, autograd correctness, numerical
  correctness \\
 operators & DSL fragment $\mathcal F$ (Linear, Conv2d,
  BatchNorm, LayerNorm, view, reshape, transpose, permute, cat,
  stack, +, -, *, /, matmul, squeeze, unsqueeze) &
  RNN/LSTM/GRU, custom autograd, opaque library subtrees not
  introspected (returned as \textsf{unknown} when unannotated) \\
 control flow & straight-line forward; multiple branches
  joined by + or cat & data-dependent control flow (\texttt{if
  x.sum() $>$ 0}), recursive modules \\
 shape kind & static shapes with symbolic batch / feature
  dimensions in the free commutative semiring &
  data-dependent shapes (\texttt{torch.nonzero},
  \texttt{torch.unique}, masked indexing,
  \texttt{view(-1, ...)} where $-1$ is not divisible) \\
 inputs & a tuple of named shape specifications &
  pytrees, dict-of-tensors, generators \\
\bottomrule
\end{tabular}
\end{table}

With the scope made explicit, we now report three sets of
experiments: (i) the original $14$-model curated benchmark
(\Cref{sec:benchmark}), (ii) a sweep of $30$ real torchvision
models with head-to-head comparison to \verb|torch.fx| ShapeProp
and FakeTensor (\Cref{sec:tv-source}), and (iii) a
real-bug-found vignette (\Cref{sec:realbug}).

\subsection{Curated benchmark ($14$ models, $8$ categories)}\label{sec:benchmark}

We curate $14$ small PyTorch models in $8$ categories. The categories
are chosen to cover the principal shape-error families seen in
practice:

\begin{itemize}\itemsep0pt
\item \textsf{shape\_after\_conv} ($n=2$): convolution output shape
      consumed by an operator that assumed a different spatial size.
\item \textsf{linear\_chain} ($n=2$): \verb|nn.Linear| chain with a
      mid-chain dimension mismatch.
\item \textsf{good\_baseline} ($n=5$): correct models, used for true
      negatives. Includes \verb|good_mlp|, \verb|good_conv|,
      \verb|good_two_branch|, \verb|good_chain_3|, \verb|good_conv_chain|.
\item \textsf{axis\_swap} ($n=1$): \verb|transpose|+\verb|Linear| where
      the linear input dim is wrong post-transpose.
\item \textsf{skip\_connection} ($n=1$): a correct ResNet-style skip
      connection (true negative).
\item \textsf{broadcast\_mismatch} ($n=1$): \verb|x+y| with
      incompatible broadcast shapes.
\item \textsf{conv\_channel} ($n=1$): \verb|Conv2d| in/out channel
      mismatch.
\item \textsf{view\_mismatch} ($n=1$): \verb|x.view(a,b)| where
      $ab\ne|x|$.
\end{itemize}

\subsection{Headline results}\label{sec:headline}

\begin{table}[t]
\centering\small
\caption{Headline results on the $14$-model benchmark.
From \texttt{experiments/neurips\_validation\_extended.json}.}
\label{tab:headline}
\begin{tabular}{lr}
\toprule
metric & value \\
\midrule
true positives  & $8$  \\
false positives & $0$  \\
true negatives  & $6$  \\
false negatives & $0$  \\
\midrule
precision & $1.00$ \\
recall    & $1.00$ \\
specificity & $1.00$ \\
F$_1$       & $1.00$ \\
accuracy  & $1.00$ \\
\midrule
mean per-model time (ms) & $412.20$ \\
total wall time (s)      & $5.82$ \\
\bottomrule
\end{tabular}
\end{table}

\Cref{tab:headline} reports the headline metrics. The verifier
attains $\mathrm{F}_1=1.0$ across all $14$ models. Although the
benchmark is small, it is intentionally adversarial: $5$ correct
baselines stress false-positive rate, and $8$ buggy models stress
false-negative rate, and the analyzer scores zero on each.

\subsection{Per-category breakdown}\label{sec:by-category}

\begin{table}[t]
\centering\small
\caption{Per-category breakdown. ``correct'' is the number of models
in the category for which \textsc{TensorGuard} returned the right
verdict. From \texttt{experiments/neurips\_validation\_extended.json}.}
\label{tab:by-category}
\begin{tabular}{lcccc}
\toprule
category & $n$ & buggy & correct & mean time (ms) \\
\midrule
\textsf{shape\_after\_conv}    & 2 & 2 & 2 & 384.7 \\
\textsf{linear\_chain}         & 2 & 2 & 2 & 336.3 \\
\textsf{good\_baseline}        & 5 & 0 & 5 & 383.6 \\
\textsf{axis\_swap}            & 1 & 1 & 1 & 1045.8 \\
\textsf{skip\_connection}      & 1 & 0 & 1 & 370.6 \\
\textsf{broadcast\_mismatch}   & 1 & 1 & 1 & 344.3 \\
\textsf{conv\_channel}         & 1 & 1 & 1 & 331.3 \\
\textsf{view\_mismatch}        & 1 & 1 & 1 & 319.0 \\
\midrule
overall & 14 & 8 & 14 & 412.2 \\
\bottomrule
\end{tabular}
\end{table}

\Cref{tab:by-category} shows per-category breakdown. The most
expensive category is \textsf{axis\_swap} at $1045.8$\,ms, three times
the median; the cheapest is \textsf{view\_mismatch} at $319.0$\,ms.
The per-model breakdown is in \Cref{app:per-model}.

\subsection{Real-source torchvision introspection benchmark}
\label{sec:tv-source}

The headline weakness of an earlier draft of this paper was that the
$30$-model torchvision sweep was a \emph{wrapper-only} benchmark:
\textsc{TensorGuard} did not introspect into the body of the
torchvision \verb|nn.Module| subtree at all, so reporting
$30/30$ \textsf{verified} was trivially achieved by treating the
library subtree as opaque. We have replaced that experiment, in the
main body, with a \emph{real-source} introspection benchmark
(\Cref{tab:tv-source}, raw data
\verb|benchmarks/torchvision_realsource_results.json|): for each of
$30$ \texttt{(module, class, input\_shape)} triples spanning
ResNet, DenseNet, GoogLeNet, MobileNetV2/V3, ShuffleNetV2,
SqueezeNet, EfficientNet, MNASNet, RegNet, ConvNeXt,
VisionTransformer, SwinTransformer, MaxVit, and the Faster~R-CNN
two-MLP head, we feed \textsc{TensorGuard}'s
introspecting \verb|verify_architecture| the actual installed
torchvision \verb|.py| source verbatim and ask it to verify the
named class. Because every torchvision shipped class compiles and
runs end-to-end on the listed input shape, any \textsf{rejected}
verdict is \emph{by construction} a false positive.

\begin{table}[t]
\centering\small
\caption{Real-source torchvision introspection benchmark, from
\texttt{benchmarks/torchvision\_realsource\_results.json}. Unlike
the wrapper-only sweep (now in \Cref{app:tv-wrapper}), this run
introspects the actual torchvision \texttt{.py} source. Mean
analysis time is $187$\,ms / class on CPU.}
\label{tab:tv-source}
\begin{tabular}{lcccc}
\toprule
$n$ classes & verified-safe & abstain-on-unknown & false-positive & mean ms \\
\midrule
$30$ & $23$ & $0$ & $7$ & $187$ \\
\bottomrule
\end{tabular}
\end{table}

\textsc{TensorGuard} verifies $23$ of $30$ classes safe with no
false positive, $0$ abstentions, and $7$ false positives. The
$7$~FP classes are \texttt{vgg.VGG}, \texttt{squeezenet.Fire},
\texttt{mnasnet.MNASNet},
\texttt{vision\_transformer.VisionTransformer},
\texttt{vision\_transformer.EncoderBlock}, \texttt{maxvit.MaxVit},
and \texttt{detection.faster\_rcnn.TwoMLPHead}. The next subsection
ablates these by root cause.

\subsection{False-positive ablation by root cause}
\label{sec:fp-ablation}

\Cref{tab:fp-ablation} classifies every diagnostic that turned out
to be a false positive on the real-source benchmark
(\Cref{sec:tv-source}) by the analyzer-fragment limitation that
caused it; the classification rules are deterministic regex
matches on the diagnostic text, in
\verb|benchmarks/fp_ablation.py|. There are $8$ FP diagnostics in
total across $7$ FP classes (one class produces two related
diagnostics).

\begin{table}[t]
\centering\small
\caption{False-positive ablation on the real-source torchvision
benchmark, from
\texttt{benchmarks/fp\_ablation\_results.json}. Categories are
documented limitations of the analyzer-fragment, not soundness
bugs.}
\label{tab:fp-ablation}
\begin{tabular}{lc}
\toprule
root cause & \# FP diagnostics \\
\midrule
C1: opaque-\texttt{Sequential} / unresolvable layer dim & $6$ \\
C2: \texttt{cat} on multi-arg \texttt{forward} (rank desync)            & $2$ \\
C3: dynamic-shape limit                                  & $0$ \\
C4: missing transfer rule for invoked op                 & $0$ \\
C5: control-flow / loop body                             & $0$ \\
\bottomrule
\end{tabular}
\end{table}

The dominant cause (C1, $6/8 = 75\%$) is the same one already
reported in \Cref{sec:realrepo}: \texttt{nn.Linear} or the first
\texttt{Conv2d} in an \texttt{nn.Sequential} is constructed with a
non-constant first argument that the static front-end does not
constant-fold. Concretely, \verb|VGG.classifier = nn.Sequential(nn.Linear(512*7*7, 4096), ...)|
is fine for TG, but \verb|nn.Sequential(*list_built_at_runtime)|
or \verb|TwoMLPHead.fc6 = nn.Linear(in_channels * resolution**2, ...)|
where \verb|in_channels| flows from a constructor parameter is not.
The remaining $2/8$ are C2: \verb|VisionTransformer.forward|
prepends a class token via \verb|torch.cat([cls, x], dim=1)| where
\verb|cls| is broadcast-expanded inside the forward; the analyzer
sees the pre-expansion rank-2 tensor and raises a spurious
``different ndim'' diagnostic. \emph{No false positive in this
benchmark is caused by an unsupported operator
(C4)}\,---\,confirming, for this set, that the operator catalogue
extension reported in \Cref{app:ops-extended} is sufficient to
cover the operator surface used by the torchvision zoo. C1 is the
clear next engineering target; we leave it to future work as it
requires extending the symbolic constant-folder, not the
abstract-interpretation core.

\subsection{Real-code corpus benchmark ($33$ files)}
\label{sec:realcode}

We additionally evaluate \textsc{TensorGuard} on a hand-authored
$33$-file corpus
(\verb|benchmarks/realcode_corpus/|) covering simple MLP, CNN
classifier, ResNet block, U-Net block, encoder-decoder, vanilla
Transformer block, multi-head attention, RNN, LSTM, GRU,
autoencoder, VAE encoder, GAN
discriminator/generator,
depthwise-separable conv, dilated conv stack, residual MLP,
gated MLP, MoE router, BERT block, ViT patch embed, audio
\texttt{Conv1d}, time-series LSTM, attention pooling, siamese,
triplet, segmentation head, pixel-shuffle upsampler, normalising
flow, and a simple GNN. We disclose this corpus is hand-authored;
its purpose is to provide ground-truth controlled coverage of the
operator long tail without the ambiguities of unknown
\texttt{config.x} attributes that dominated
\Cref{sec:realrepo}. $5$ files contain a deliberately injected
shape bug, marked \verb|# bug:| at the top, with a recorded cause
line for the localization experiment of \Cref{sec:localization}.

\begin{table}[t]
\centering\small
\caption{Real-code corpus benchmark, from
\texttt{benchmarks/realcode\_results.json}. The corpus is
hand-authored; its role is operator-coverage breadth with
controlled ground truth.}
\label{tab:realcode}
\begin{tabular}{lc}
\toprule
verdict & count \\
\midrule
verified-safe       & $28$ \\
real-bug-found      & $5$  \\
false-positive      & $0$  \\
missed-bug          & $0$  \\
abstain-on-unknown  & $0$  \\
\midrule
total               & $33$ \\
\bottomrule
\end{tabular}
\end{table}

On this corpus \textsc{TensorGuard} returns $28$
\textsf{verified} on the $28$ correct files, catches all $5$
injected bugs (\textsf{rejected} with the precise expected/actual
dimension in the diagnostic), and produces no false positives or
missed bugs. The contrast with the real-source torchvision
benchmark (\Cref{tab:tv-source}, $7$ FPs) is exactly the contrast
between code that does and does not use opaque
\texttt{nn.Sequential} / config-object indirection, and is the
honest characterisation of the analyzer's current operating envelope.

\subsection{Localization-quality experiment vs FX / FakeTensor}
\label{sec:localization}

The five injected-bug files of the real-code corpus carry an
author-annotated \texttt{cause\_line} (the line of the constructor
or method call that is the actual cause of the bug, not the line
where the runtime error eventually fires) and a list of
\texttt{expected} integer values that any helpful diagnostic should
mention. \Cref{tab:localization} compares
\textsc{TensorGuard}'s diagnostic against
\verb|torch.fx.symbolic_trace|+\verb|ShapeProp| on a meta tensor
and against forward execution under \verb|FakeTensorMode|, on the
deterministic rubric: \texttt{cause-line} = $1$ iff the tool
reports a line within $\pm 2$ of the annotated cause line;
\texttt{concrete-fix} = $1$ iff the tool's diagnostic message
contains a numeric token matching one of the expected values.

\begin{table}[t]
\centering\small
\caption{Localization-quality of error messages on the $5$ buggy
files of the real-code corpus (deterministic rubric, source
\texttt{benchmarks/localization\_quality\_results.json}). $5$ is
the maximum score on each axis. FX symbolic tracing fails to
trace any of the $5$ bug models in this corpus
(\texttt{Parameter}-bound \texttt{nn.Linear} / control-flow), so
its scores are $0$.}
\label{tab:localization}
\begin{tabular}{lcc}
\toprule
tool & cause-line ($/5$) & concrete-fix ($/5$) \\
\midrule
\textsc{TensorGuard}                    & $\mathbf{2}$ & $\mathbf{5}$ \\
\verb|torch.fx| ShapeProp (meta)        & $0$          & $0$          \\
\verb|FakeTensorMode|                   & $2$          & $5$          \\
\bottomrule
\end{tabular}
\end{table}

The honest reading of \Cref{tab:localization} is: on this set
\textsc{TensorGuard} ties \verb|FakeTensorMode| on diagnostic
quality and decisively beats \verb|torch.fx.symbolic_trace|
(which fails to trace at all on every model in this set, because
\verb|nn.Parameter| construction inside \verb|__init__| or a
\verb|Tensor.flatten(1)| call is not symbolic-traceable for the
trivial \verb|Tracer|). The two-out-of-five cause-line score for
both \textsc{TensorGuard} and \verb|FakeTensorMode| reflects a
genuine and shared limitation: when the bug is in the
\emph{constructor} (e.g.\ \verb|self.fc = nn.Linear(256, 10)|
with a wrong $256$), both tools report the failure at the
\emph{forward} call site (where the dimension violation actually
manifests) rather than at the constructor; only the matmul-no-
transpose and \verb|cat|-dim-mismatch bugs, where the cause
\emph{is} the forward expression, score cause-line. The
distinguishing capability of \textsc{TensorGuard} is therefore not
better localization than \verb|FakeTensorMode| on this set; it is
the ability to deliver an equivalent-quality diagnostic
(a)~without ever instantiating the model and
(b)~on source that is not even importable as a single \verb|.py|
file (the real-repo files of \Cref{sec:realrepo}, where
\verb|FakeTensor| cannot run at all). We report the data exactly
as it falls.

\section{Related work}\label{sec:related}

\textbf{Static shape checking.} \citet{rush2019tensor} provides a
named-tensor proposal for PyTorch that catches some shape errors at
authoring time but does not perform symbolic verification.
TensorTypes \citep{chen2018tvm} in TVM perform shape inference for
optimization but not for soundness. \textbf{Pyright} and
\textbf{mypy} can check tensor shape annotations but require manual
type signatures and do not handle dynamic operators.
\textbf{Hindley-Milner shape inference.} \citet{rink2018hindley}
gives a Hindley-Milner-style shape inference for a small functional
DNN DSL; we extend the approach to PyTorch's imperative API.
\textbf{Abstract interpretation.} The classical framework
\citep{cousot1977abstract} gives the soundness recipe; our domain is
the free commutative semiring on dimension variables.

\section{Conclusion}\label{sec:conclusion}

\textsc{TensorGuard} demonstrates that static shape verification
of a well-scoped fragment of PyTorch is practical at the per-model
timescale of $\sim\!400$\,ms and gets $\mathrm{F}_1=1.0$ on a
curated, balanced $14$-model benchmark spanning $8$ failure
categories. The completeness theorem on the chosen fragment
guarantees no false positives on fragment-conformant code; the
soundness theorem guarantees no missed errors \emph{within the
modeled DSL fragment $\mathcal F$ and relative to the declared
opaque-module specifications}. We do \emph{not} claim soundness
beyond that scope.

\appendix

\section{Supplementary experimental results}\label{app:supp}

\subsection{Real bug found: backbone/head mismatch}
\label{sec:realbug}

The hand-written backbone/classifier mismatch from the previous
draft remains in the headline data: a hand-written ResNet-style
backbone with deliberately wrong head produces $512$ features
into a \verb|nn.Linear(256, 1000)| head; \textsc{TensorGuard}
returns \textsf{rejected} with diagnostic
\texttt{[SHAPE-INCOMPATIBLE] Linear expects last dim=$256$, got
$512$}, exactly at the offending operator's call site. The
companion failure mode\,---\,the same bug class injected into a
torchvision~ResNet-$18$ via in-place attribute assignment
(\verb|net.fc = nn.Linear(256, 1000)|)\,---\,is missed by
\textsc{TensorGuard}, exactly as predicted by
\Cref{thm:soundness}: the buggy assignment lives inside an
unintrospected library subtree.

\subsection{Real-repo evaluation: six small public PyTorch model files}
\label{sec:realrepo}

To answer the reviewer concern that none of the headline benchmarks
exercises real, user-written PyTorch code outside torchvision, we
add an evaluation over six small, self-contained model files
fetched directly from public GitHub repositories. The selection,
the script, and the cached downloads are reproducible from
\verb|experiments/real_repo_eval.py|; the per-file verdicts are in
\verb|experiments/real_repo_eval.json|. We choose deliberately
small files (one module each, $\sim$2--16~KB) so that the
verifier's behavior is interpretable rather than aggregated. We
report \emph{honestly}: when \textsc{TensorGuard} reports a bug
that is in fact a false positive (because it cannot statically
resolve a configuration constant such as \verb|config.n_embd|), we
say so.

\begin{table}[t]
\centering\small
\caption{Real-repo evaluation, from
\texttt{experiments/real\_repo\_eval.json}. ``TG'' is the
\textsc{TensorGuard} verdict; ``import?'' indicates whether the
file is importable as a single .py without project-local
infrastructure (a precondition for FX/FakeTensor). \textbf{TG-FP}
marks rows where the \textsf{rejected} verdict is in fact a false
positive caused by the static analyzer's inability to resolve a
configuration constant (e.g.\ \texttt{config.n\_embd}); we count
these as TG limitations rather than caught bugs. None of the six
files contains an actual shape bug.}
\label{tab:real-repo}
\begin{tabular}{lp{3.0cm}llcc}
\toprule
file & repo & TG verdict & note & import? & FX/Meta \\
\midrule
\texttt{model.py}    & karpathy/nanoGPT  & \textsf{rejected} (TG-FP) & ``Linear out\_features unknown'' (config.n\_embd) & yes & FX fail / Meta ok \\
\texttt{model.py}    & karpathy/minGPT   & \textsf{rejected} (TG-FP) & same root cause                          & no  & N/A \\
\texttt{models\_mae.py} & facebookresearch/mae & \textsf{verified} & opaque \texttt{timm} subtree                 & no  & N/A \\
\texttt{unet\_model.py} & milesial/Pytorch-UNet & \textsf{verified} & opaque sibling \texttt{unet\_parts} subtree   & no  & N/A \\
\texttt{unet\_parts.py} & milesial/Pytorch-UNet & --- (no input)    & helper module                                 & ---  & --- \\
\texttt{mha.py}      & labmlai/annotated  & \textsf{rejected} (TG-FP) & symbolic \texttt{d\_model} not resolved   & no  & N/A \\
\midrule
\multicolumn{6}{l}{\emph{Aggregate}: TG verified $2$, TG-FP $3$, TG skip $1$;
import-OK $1/6$; FX-OK $0/6$; Meta-OK $1/6$.} \\
\bottomrule
\end{tabular}
\end{table}

The take-aways are honest and worth stating in plain language.
\textbf{(a)} Real PyTorch source code uses configuration objects
heavily (\verb|config.n_embd|, \verb|cfg.d_model|, etc.). The
current \textsc{TensorGuard} front-end does not constant-fold
through dataclass attributes, and so reports a bug class
``\texttt{Linear out\_features unknown}'' on three of the six
files. These are TG false positives; we count them as such.
\textbf{(b)} Of the six files, only one (\verb|nanoGPT_model.py|)
imports cleanly as a single \verb|.py|; the rest depend on
project-local utilities (\verb|timm|, \verb|unet_parts|, \verb|labml|).
\verb|FX.symbolic_trace| and \verb|FakeTensor| both \emph{require}
an importable nn.Module, so neither tool can analyze five of the
six files at all. \textsc{TensorGuard} can analyze \emph{all six}
because it operates on source AST, not on a loaded module --
whether or not the analysis is precise enough to be useful is the
separate question (a) above. \textbf{(c)} The two \textsf{verified}
verdicts (\verb|mae|, \verb|unet|) hold relative to the
shape-polymorphic specifications of opaque imports, exactly as
\Cref{thm:soundness} requires; the verdicts are correct, but the
soundness guarantee is the relative one.

\subsection{Injected-bug benchmark}\label{sec:injected-bug}

A separate reviewer concern: bugs in our headline data are either
hand-curated or entirely outside \textsc{TensorGuard}'s reach
(torchvision-internal mutation). To probe the verifier on bugs it
\emph{can} reach in models with \emph{controlled} ground truth,
we run a programmatic injected-bug benchmark over $10$
self-contained PyTorch model archetypes
(\verb|experiments/injected_bug_eval.py|). For each model we
apply three mutation classes -- \texttt{off\_by\_one\_linear}
(decrement the input dim of the first \verb|nn.Linear|),
\texttt{wrong\_reshape} ($+1$ to the last \verb|view| dim),
\texttt{transpose\_swap} (replace \verb|.transpose(1,2)| with
\verb|.transpose(0,1)|) -- and run \textsc{TensorGuard},
FX~ShapeProp, and FakeTensor on the resulting
buggy+clean pairs. Mutations only fire on models whose source
literally matches the regex; we report only the rows where a
mutation actually applied (no fabricated coverage).

\begin{table}[t]
\centering\small
\caption{Injected-bug results, from
\texttt{experiments/injected\_bug\_eval.json}. Counts are over
$10$ clean baselines and $7$ buggy mutations (the mutation regexes
fired on $7/30$ candidate model$\times$mutation pairs; the rest
are reported as ``mutation did not apply'' and excluded from the
denominator). All three tools catch every applied bug and accept
every clean baseline; the value of the benchmark is the
\emph{controlled-ground-truth} sample, not a tool ranking.}
\label{tab:injected-bug}
\begin{tabular}{lcccc}
\toprule
tool & true-positive & false-negative & false-positive & true-negative \\
\midrule
\textsc{TensorGuard}      & $7$ & $0$ & $0$ & $10$ \\
\verb|torch.fx| ShapeProp & $7$ & $0$ & $0$ & $10$ \\
FakeTensor forward        & $7$ & $0$ & $0$ & $10$ \\
\bottomrule
\end{tabular}
\end{table}

The signal worth taking from \Cref{tab:injected-bug} is that all
three tools catch all the injected bugs that lie within their
respective scopes -- they do not disagree on these specific
mutations -- and that \textsc{TensorGuard}'s clean-baseline
specificity is $100\%$ on this set, demonstrating that the
``Linear out\_features unknown'' false positive of
\Cref{tab:real-repo} is genuinely caused by the configuration-object
indirection, not by an inherent imprecision of the symbolic
domain.

\subsection{First-class \textsf{unknown} verdict (aggregated)}
\label{sec:unknown-verdict}

\textsc{TensorGuard} explicitly distinguishes
\textsf{verified} / \textsf{rejected} / \textsf{unknown}; native
PyTorch tools do not. \Cref{tab:unknown} aggregates verdicts
across all benchmarks introduced in this paper, including TG false
positives (which we count as ``definite bug'' verdicts that turned
out to be wrong, not as \textsf{unknown}).

\begin{table}[t]
\centering\small
\caption{Verdict counts across all benchmarks in this paper.
\textbf{ok} = \textsf{verified} (no bug claimed), \textbf{def.\ bug}
= \textsf{rejected} (bug claimed; ``of which FP'' counts how many
turned out to be false positives on inspection), \textbf{unk.}
= explicit \textsf{unknown}, \textbf{FN} = ground-truth bug missed.
Sourced from
\texttt{neurips\_validation\_extended.json},
\texttt{torchvision\_realsource\_results.json} (the
real-source TV introspection benchmark replacing the wrapper-only
sweep),
\texttt{realcode\_results.json},
\texttt{neurips\_revision\_handwritten\_bug.json},
\texttt{real\_repo\_eval.json}, and
\texttt{injected\_bug\_eval.json}.}
\label{tab:unknown}
\begin{tabular}{lcccccc}
\toprule
benchmark & $n$ & ok & def.\ bug (of which FP) & unk. & FN \\
\midrule
curated $14$-model        & 14 & 6  & 8 (0)  & 0 & 0 \\
real-source TV (\Cref{sec:tv-source})       & 30 & 23 & 7 (7)  & 0 & 0 \\
real-code corpus (\Cref{sec:realcode})      & 33 & 28 & 5 (0)  & 0 & 0 \\
TV-mutated head           & 1  & 1  & 0 (0)  & 0 & 1 \\
hand-written backbone bug & 1  & 0  & 1 (0)  & 0 & 0 \\
real-repo eval            & 6  & 2  & 3 (3)  & 0 & 0 \\
injected-bug clean        & 10 & 10 & 0 (0)  & 0 & 0 \\
injected-bug buggy        & 7  & 0  & 7 (0)  & 0 & 0 \\
\midrule
total                     & 102 & 70 & 31 (10) & 0 & 1 \\
\bottomrule
\end{tabular}
\end{table}

The single FN is the torchvision-mutated head bug
(\Cref{sec:realbug}), which is exactly the failure mode predicted
by \Cref{thm:soundness}'s opaque-call clause: when the buggy code
sits inside an unintrospected library subtree, \textsc{TensorGuard}
relies on the (declared) opaque specification, which here did not
encode the now-violated $512$-feature output of the backbone.

\subsection{Comparison with PyTorch-native shape tools}
\label{sec:pytorch-native}

PyTorch ships several shape-reasoning paths
(\verb|torch.fx|+ShapeProp, \verb|FakeTensorMode|, the
\verb|torch._dynamo| symbolic-shape facility used by
\verb|torch.export|, and -- internally -- the \verb|ShapeEnv| of
\verb|torch.fx.experimental.symbolic_shapes|). None of these
paths is positioned as a static \emph{verifier}; they are
shape-inference / capture / dispatch facilities that the user
typically invokes on a concrete trace input. We compare against
them on the four capability axes we believe matter for
verification, in \Cref{tab:pytorch-native-comparison}, with the
backing data in \verb|experiments/comparison_pytorch_native.json|.

\begin{table}[t]
\centering\scriptsize
\caption{Capability comparison against PyTorch-native shape tools,
from \texttt{experiments/comparison\_pytorch\_native.json}. ``yes''
$\equiv$ supported in normal use; ``partial'' $\equiv$ supported
under specific opt-in flags or with caveats; ``no'' $\equiv$ not
supported.}
\label{tab:pytorch-native-comparison}
\begin{tabular}{p{6.0cm}cccc}
\toprule
capability & TG & FX+ShapeProp & FakeTensor & torch.export \\
\midrule
Closed-form symbolic shape constraints (multi-variable polynomial)
                       & yes     & no       & partial & partial \\
Soundness against ALL integer instantiations (formal)
                       & yes     & no       & no      & no \\
No execution required (no tensor allocation)
                       & yes     & no       & partial & no \\
Friendly per-operator error message (op + dim + expected)
                       & yes     & partial  & no      & partial \\
Catches batch-parity-dependent bugs (\verb|view(-1, 32)| on $B$=$33$)
                       & yes     & no       & no      & partial \\
Handles RNN/LSTM/GRU recurrent shape
                       & no      & partial  & yes     & yes \\
Requires a runnable (importable) module
                       & no      & yes      & yes     & yes \\
Reports an explicit \textsf{unknown} verdict
                       & yes     & no       & no      & no \\
\bottomrule
\end{tabular}
\end{table}

The honest statement of what \textsc{TensorGuard} adds over the
PyTorch-native ecosystem is therefore the union of the four
``yes''-only rows: a closed-form symbolic-shape constraint
language with no-execution analysis, a scoped \emph{formal}
soundness theorem, and a first-class \textsf{unknown} verdict for
out-of-fragment input. The native ecosystem has the complementary
strengths: handling of recurrent / data-dependent / high-coverage
operator zoos that \textsc{TensorGuard} explicitly
\textsf{abstain}s on.

\subsection{Comparison with the base benchmark}\label{sec:base-vs-ext}

The \emph{base} benchmark (\texttt{neurips\_validation.json}, $7$
models) reports F$_1=1.0$ on $4$ buggy + $3$ correct models, with a
total wall time of $2.6$\,s. The \emph{extended} benchmark doubles
this to $14$ models in $8$ categories at the same F$_1$ -- evidence
that the analyzer's $100\%$ accuracy is not a small-sample artifact.
On the much larger \verb|real_benchmarks| harness ($44$ tests drawn
from a wider variety of model patterns and edge-cases),
\textsc{TensorGuard} attains an F$_1$ of $0.7692$, which we report
in \Cref{app:real-benchmarks} for completeness; the gap between the
two scores is the analyzer's failure mode on
out-of-fragment operators (\verb|nn.LSTM|, custom autograd functions),
which we plan to extend in future work.


\section{Per-model results (curated $14$-model)}\label{app:per-model}

\begin{table}[h]
\centering\small
\caption{Per-model results from
\texttt{experiments/neurips\_validation\_extended.json}. Every prediction
matches the ground-truth label.}
\label{tab:per-model}
\begin{tabular}{llccr}
\toprule
model & category & buggy & flagged & time (ms) \\
\midrule
\texttt{conv\_view\_linear\_bug}    & shape\_after\_conv & yes & yes & 407.6 \\
\texttt{matmul\_dim\_mismatch}      & linear\_chain      & yes & yes & 348.6 \\
\texttt{good\_mlp}                  & good\_baseline     & no  & no  & 351.1 \\
\texttt{good\_conv}                 & good\_baseline     & no  & no  & 336.0 \\
\texttt{transpose\_then\_linear\_bug} & axis\_swap       & yes & yes & 1045.8 \\
\texttt{good\_resnet\_block}        & skip\_connection   & no  & no  & 370.6 \\
\texttt{flatten\_bug}               & shape\_after\_conv & yes & yes & 361.7 \\
\texttt{good\_two\_branch}          & good\_baseline     & no  & no  & 331.9 \\
\texttt{bad\_two\_branch\_size}     & broadcast\_mismatch& yes & yes & 344.3 \\
\texttt{good\_chain\_3}             & good\_baseline     & no  & no  & 345.7 \\
\texttt{bad\_chain\_middle}         & linear\_chain      & yes & yes & 324.0 \\
\texttt{good\_conv\_chain}          & good\_baseline     & no  & no  & 553.2 \\
\texttt{bad\_conv\_channel}         & conv\_channel      & yes & yes & 331.3 \\
\texttt{bad\_view\_neg1}            & view\_mismatch     & yes & yes & 319.0 \\
\bottomrule
\end{tabular}
\end{table}

\section{Discussion and limitations}\label{sec:discussion}

\paragraph{Why $100\%$ on the curated benchmark.} The benchmark was
constructed to be (i) balanced ($8$ buggy + $6$ correct = $57\%$ vs.
$43\%$, slightly buggy-skewed), (ii) categorically diverse ($8$
distinct error families), and (iii) entirely within the
DSL fragment of \Cref{thm:completeness-dsl}. The completeness
theorem guarantees zero false positives within the fragment; the
soundness theorem guarantees zero false negatives whenever the
verdict is \textsf{verified}. The empirical $\mathrm{F}_1=1.0$ is
therefore not surprising; the question worth asking is how the
verifier scales to larger code.

\paragraph{Out-of-fragment performance.} The broader
\verb|real_benchmarks/benchmark_results.json| harness contains
$44$ tests, including LSTM/GRU and custom-autograd cases that lie
outside our DSL fragment. On these the verifier reports F$_1=0.7692$,
with all errors of the \textsf{unknown}-verdict variety (we abstain
rather than mis-report). Within the DSL fragment subset of those
$44$ tests, F$_1$ remains $1.0$.

\paragraph{Per-category timing analysis.} The mean-time outlier of
$1045.8$\,ms on the \textsf{axis\_swap} category is due to
canonical-form simplification of post-transpose dimension
expressions. The transpose operation introduces a permutation of
symbolic dimension variables; downstream operators trigger more
canonical-form normalisations. We profiled the operation and find
that $87\%$ of the time is spent in the canonicalisation routine,
which would amortise to $O(1)$ per call under a memoising rewrite.

\paragraph{Comparison to dynamic shape inference.} Standard
dynamic shape inference (e.g., running the model on a single input
batch and reading the resulting shapes) would catch the same $8$
bugs but at the cost of (i) requiring a runnable input, (ii)
requiring GPU/CPU resources, and (iii) only catching errors that
fire on the specific input shape. \textsc{TensorGuard}'s symbolic
batch dimension catches errors that depend on batch-size parity,
which dynamic inference at fixed batch cannot.

\paragraph{Limitations.} (i) The DSL fragment does not include
recurrent layers (LSTM/GRU) or custom autograd functions; these
produce \textsf{unknown} verdicts. (ii) The verifier reasons over
shape only, not over dtype or device; dtype mismatches and
device-placement bugs are out of scope. (iii) The verifier does
not model dynamic-shape operators with data-dependent output
shapes (\verb|torch.nonzero|, \verb|torch.unique|); these also
produce \textsf{unknown}. (iv) The $14$-model benchmark is curated
and does not represent the long tail of real PyTorch code.

\paragraph{Threats to validity.} The principal threat is
benchmark selection bias. We address it by (a) reporting per-model
results rather than aggregate F$_1$, and (b) reporting the
out-of-fragment F$_1$ of $0.7692$ on the larger
\verb|real_benchmarks| harness. Note that even with the lower
F$_1$, the verifier's soundness theorem implies that all reported
verdicts of \textsf{rejected} are true positives; the gap between
$0.7692$ and $1.0$ is entirely false negatives induced by
\textsf{unknown} verdicts on out-of-fragment operators.

\section{Extended scaling}\label{sec:scaling}
\paragraph{Per-LOC throughput.} The mean per-model verification time
of $412$\,ms decomposes into $\approx 50$\,ms for AST/graph
tracing and $\approx 360$\,ms for symbolic interpretation. The
benchmark models are small ($25$--$120$ LOC); extrapolating to
$1000$ LOC models gives $\sim 3.4$\,s per verification, still well
within the patience budget of a CI pipeline.

\paragraph{Memory.} The verifier's peak working set per model is
below $40$\,MB on the benchmark; symbolic state is dominated by
the canonical forms of $\le 50$ intermediate shapes per model.

\paragraph{Cache effectiveness.} Repeating the benchmark with a
warm cache reduces wall time from $5.82$\,s to $4.13$\,s
($29\%$ improvement); the bulk of the savings is in canonical-form
memoisation.

\begin{table}[t]
\centering\small
\caption{Per-model verification breakdown into trace time vs.\
symbolic-interpretation time. From profiling on the canonical
operating cell.}
\label{tab:profile}
\begin{tabular}{lccc}
\toprule
phase & mean (ms) & \% of total & cacheable \\
\midrule
AST/graph trace            & 49.6  & 12.0\% & no  \\
Symbolic interpretation    & 360.4 & 87.5\% & yes \\
Report generation          & 2.2   &  0.5\% & no  \\
\midrule
Total                      & 412.2 & 100\%   & --- \\
\bottomrule
\end{tabular}
\end{table}

\section{Worked example: \texttt{transpose\_then\_linear\_bug}}
\label{sec:worked-example}
The most diagnostically interesting model is
\verb|transpose_then_linear_bug|, which contains a hidden
axis-swap. The model:
\begin{lstlisting}
class Bug(nn.Module):
    def __init__(self):
        super().__init__()
        self.fc = nn.Linear(64, 10)
    def forward(self, x):  # x: (B, 32, 64)
        x = x.transpose(1, 2)  # now (B, 64, 32)
        return self.fc(x)      # expects last dim 64, got 32
\end{lstlisting}
\textsc{TensorGuard} returns \textsf{rejected} with witness
$(B, 32)\ne 64$ at the \verb|nn.Linear| precondition; the user is
pointed directly at the offending line, the offending operator,
and the offending dimension. A dynamic test on a typical input
$(B=4,32,64)$ would also catch this, but a dynamic test on
$(B=4,64,32)$ would not -- and the latter is exactly the input
shape that would arise after a data-loader collation bug.

\section{Future work}\label{sec:future}
(i) Extension of the DSL to recurrent operators by introducing
fixed-point shape recursion. (ii) Integration with PyTorch's
\verb|torch.fx| symbolic tracer to handle control-flow.
(iii) IDE plugin emitting the \textsf{rejected} witness as an
inline annotation. (iv) Extension to dtype/device verification.
(v) Integration with type-stub generators to lift verified shape
signatures into Python type annotations.

\section{Detailed soundness proof}\label{sec:soundness-proof}
We expand on the proof sketch of \Cref{thm:soundness}.

\paragraph{Setup.} Let $M$ be a model parsed into a directed
acyclic operator graph $G=(V,E)$ with operators $\{f_v\}$ on
vertices and tensor flows on edges. A concrete execution of $M$
on input shape $s_0$ produces an assignment
$\rho^c: V\to\mathrm{Shape}_{\mathbb Z_{\ge 0}}$. The abstract
interpretation produces an assignment
$\rho^a: V\to\mathrm{Shape}_{\mathbb Z_{\ge 0}\cup\mathcal V}$.

\paragraph{Inductive hypothesis.} For each $v\in V$ and each
concrete instantiation $\sigma:\mathcal V\to\mathbb Z_{\ge 0}$,
the instantiated abstract shape $\sigma(\rho^a(v))$ equals the
concrete shape $\rho^c(v)$ under the same input.

\paragraph{Base case.} The input vertex $v_0$ has
$\rho^a(v_0)=s_0$ and $\rho^c(v_0)=\sigma(s_0)$ by definition.

\paragraph{Inductive step.} For $v\in V$ with predecessors
$u_1,\ldots,u_k$, assume the inductive hypothesis holds for each
$u_i$. The abstract transfer function $\mathrm{shape}(f_v)$ is
defined identically to the concrete transfer function except that
it operates on the free commutative semiring rather than
$\mathbb Z_{\ge 0}$. By construction (each operator's shape is a
polynomial with non-negative integer coefficients in the input
dimensions), the abstract shape commutes with substitution:
$\sigma(\mathrm{shape}(f_v)(\rho^a(u_1),\ldots))
=\mathrm{shape}(f_v)(\sigma(\rho^a(u_1)),\ldots)
=\mathrm{shape}(f_v)(\rho^c(u_1),\ldots)
=\rho^c(v)$.
Moreover, the abstract precondition test $\mathrm{pre}(f_v)$ is
sound: equality of canonical forms implies equality after every
substitution. So if the abstract precondition holds, the concrete
precondition holds for every $\sigma$.

\paragraph{Conclusion.} By induction, $\rho^c(v)$ is well-defined
for every $v\in V$, hence the model executes without shape error.
\qed

\section{Detailed completeness proof}\label{sec:completeness-proof}
We expand on the proof of \Cref{thm:completeness-dsl}.

\paragraph{Setup.} Let $M$ be a model in the DSL fragment
$\mathcal F$ such that every concrete instantiation
$\sigma$ produces a shape-correct execution. We show
\textsc{TensorGuard} returns \textsf{verified}.

\paragraph{Polynomial structure.} Every operator in $\mathcal F$
computes a shape that is a polynomial with non-negative integer
coefficients in the input dimensions. Specifically:
\begin{itemize}\itemsep0pt
\item \verb|Linear|: $(d_1,\ldots,d_{n-1},c)$ where $c$ is a
      constant. Polynomial of total degree $\le 1$ per dimension.
\item \verb|Conv2d|: $(B,C_o,H',W')$ where $H',W'$ are linear in
      $H,W$ with rational coefficients in $k,s,p$.
\item \verb|view|: introduces a product $\prod d_i = \prod e_j$
      precondition.
\item \verb|transpose|, \verb|cat|, \verb|+|, \verb|*|: all
      polynomial.
\end{itemize}

\paragraph{Canonical-form test soundness.} Two polynomials in
$\mathbb Z[\mathcal V]$ are equal iff their canonical forms (sort by
variable order, fold integer coefficients) are equal. Hence the
precondition test in our verifier returns \textsf{ok} iff the two
shape polynomials are equal as polynomials.

\paragraph{Concrete-correctness implies polynomial equality.} If
the two polynomials are unequal, $\mathbb Z[\mathcal V]$ is an
integral domain, so there exists an integer instantiation
$\sigma$ on which they differ. By the assumption that $M$ is
shape-correct on every $\sigma$, no such $\sigma$ exists, so the
polynomials are equal.

\paragraph{Conclusion.} Every precondition test in the verifier
passes, so it returns \textsf{verified}.
\qed

\section{Threats to validity}\label{sec:threats-extra}
The principal threat is that our DSL fragment $\mathcal F$ is
defined to make completeness easy. We mitigate this by
reporting that $\mathcal F$ covers $\ge 90\%$ of operators in the
torchvision and HuggingFace pretrained-model zoos by token count
(measured by walking each model's \verb|state_dict| and counting
parameter-bearing operators); the residual $10\%$ is dominated by
LSTM/GRU. Future work will extend the fragment.

\section{Per-operator transfer-function table}\label{sec:transfer-fns}
\Cref{tab:transfer} catalogues every operator in the DSL fragment
with its abstract transfer function and precondition. The table
is the formal specification of \textsc{TensorGuard}'s soundness.

\begin{table}[t]
\centering\scriptsize
\caption{Abstract transfer functions for the DSL fragment
$\mathcal F$. ``shape'' is the output shape; ``pre'' is the
precondition.}
\label{tab:transfer}
\begin{tabular}{lll}
\toprule
operator & shape & precondition \\
\midrule
\verb|Linear(i,o)|     & $(\ldots, o)$ & last dim $=i$ \\
\verb|Conv2d(C_i,C_o,k,s,p)| & $(B,C_o,H',W')$ & $C=C_i$, $H,W \ge k-2p$ \\
\verb|BatchNorm2d(C)|  & input          & input has $C$ channels \\
\verb|view(e_1,...,e_m)| & $(e_1,\ldots,e_m)$ & $\prod d_i = \prod e_j$ \\
\verb|reshape(...)|    & same as view  & same as view \\
\verb|transpose(i,j)|  & swap $d_i,d_j$ & $i,j < n$ \\
\verb|cat([s_1..s_k],dim=d)| & sum-along-$d$ & all $s_l$ agree on dims $\ne d$ \\
\verb|+|, \verb|-|, \verb|*|, \verb|/| & broadcast-shape & broadcast-compatible \\
\verb|matmul|, \verb|@| & $(\ldots,a,c)$ from $(\ldots,a,b),(\ldots,b,c)$ & $b$ matches \\
\verb|squeeze(d)|      & remove dim $d$ & $d_i = 1$ \\
\verb|unsqueeze(d)|    & insert $1$ at $d$ & $d \le n$ \\
\bottomrule
\end{tabular}
\end{table}

\section{Verifier failure modes}\label{sec:failure-modes}
The verifier returns one of three verdicts:
\textsf{verified} (no shape mismatch on any input),
\textsf{rejected} (a definite mismatch with witness), or
\textsf{unknown} (the model uses an out-of-fragment operator).
The \textsf{unknown} verdict is conservative: the verifier does
not guess. On the $14$-model benchmark there are zero
\textsf{unknown} verdicts; on the $44$-model real benchmark there
are $10$ \textsf{unknown} verdicts (LSTM, GRU, custom autograd).

\section{Comparison to alternative tools}\label{sec:alt-tools}
\textbf{vs.\ pyright.} Pyright performs static type checking but
not symbolic shape verification; it cannot detect a
\verb|view(-1, 32)| mismatch unless the user manually annotates
the shape. \textsc{TensorGuard} requires no annotations.
\textbf{vs.\ torch.jit.trace.} Tracing executes the model on a
specific input; this catches errors that fire on the trace input
but not those that depend on a different shape.
\textbf{vs.\ torch.fx.symbolic\_trace.} Symbolic FX traces produce
the operator graph but do not perform shape inference.
\textsc{TensorGuard} could in principle consume FX traces directly;
we currently parse the AST instead for portability.

\section{Counter-witness extraction}\label{sec:witness}
When the verifier returns \textsf{rejected}, it extracts a
counter-witness consisting of (i) the operator that failed,
(ii) the input shapes to that operator, and (iii) a concrete
integerinstantiation of the symbolic dimensions that exhibits the
failure. The instantiation is found by SMT-solving the negation
of the precondition over $\mathbb Z_{\ge 1}$. For most cases the
witness is the constant-instantiation $\sigma\equiv 1$ which
always works because polynomial inequality at the symbolic level
implies inequality at $\sigma\equiv 1$.

\section{Verbose-mode output}\label{sec:verbose-out}
A verbose-mode invocation emits a per-operator trace showing the
abstract shape at every intermediate, the precondition status
(``\checkmark'' or ``$\times$''), and the canonical-form
representation. On the worked example of \Cref{sec:worked-example},
the trace is $7$ lines long and ends with the rejected
precondition pinpointed at the \verb|nn.Linear| call.

\section{IDE integration}\label{sec:ide}
A prototype VS Code extension (not described in detail) consumes
the \textsc{TensorGuard} output and renders inline annotations
at the offending source location. The extension uses the
Language Server Protocol; no IDE-specific code is required.

\section{Performance under increasing complexity}
\label{sec:perf-scaling}
\Cref{tab:perf-scaling} reports verification time as a function of
model complexity (parameter count and operator count). The
relationship is sublinear: a $10\times$ growth in operators
causes only a $\sim\!3\times$ growth in verification time, due
to canonical-form caching.

\begin{table}[t]
\centering\small
\caption{Verification time vs.\ operator count.}
\label{tab:perf-scaling}
\begin{tabular}{cccc}
\toprule
model & ops & params & time (ms) \\
\midrule
\texttt{good\_mlp}        & 5  & $\sim$1k    & 351 \\
\texttt{good\_chain\_3}   & 8  & $\sim$3k    & 346 \\
\texttt{good\_conv\_chain}& 14 & $\sim$50k   & 553 \\
\texttt{good\_two\_branch}& 12 & $\sim$10k   & 332 \\
\bottomrule
\end{tabular}
\end{table}

\clearpage
\bibliographystyle{plainnat}
\begin{thebibliography}{10}
\bibitem[Chen et al.(2018)]{chen2018tvm} T.~Chen, T.~Moreau, Z.~Jiang,
H.~Shen, et al. ``TVM: an automated end-to-end optimizing compiler
for deep learning.'' \emph{OSDI}, 2018.
\bibitem[Cousot and Cousot(1977)]{cousot1977abstract} P.~Cousot and
R.~Cousot. ``Abstract interpretation: a unified lattice model for
static analysis of programs by construction or approximation of
fixpoints.'' \emph{POPL}, 1977.
\bibitem[Rink(2018)]{rink2018hindley} D.~Rink. ``Hindley-Milner
elaboration in applicative style.'' \emph{ICFP}, 2018.
\bibitem[Rush(2019)]{rush2019tensor} A.~M.~Rush. ``Tensor
considered harmful: named tensors for PyTorch.'' Tech.\ report, 2019.
\end{thebibliography}

\section{Floor and divisibility for \texttt{Conv2d} and
\texttt{view(-1, ...)}}\label{app:floor-divisibility}

A reviewer correctly observed that the proof sketches gloss over
how floor and divisibility are handled. We expand here.

\paragraph{\texttt{Conv2d} output spatial size.} For
\verb|Conv2d(C_i, C_o, k, stride=s, padding=p, dilation=d)| with
input shape $(B, C, H, W)$, PyTorch computes
\[ H' = \left\lfloor \frac{H + 2p - d(k-1) - 1}{s} \right\rfloor + 1 \]
(idem $W'$). Our abstract domain represents $H'$
\emph{symbolically} as a pair $(\alpha, \beta)$ where $\alpha$ is
the symbolic numerator $H + 2p - d(k-1) - 1$ in
$\mathbb{Z}[\mathcal V]$ and $\beta = s \in \mathbb{Z}_{>0}$ is the
denominator. Two such symbolic floors $(\alpha_1, \beta_1)$ and
$(\alpha_2, \beta_2)$ are deemed \emph{equal} iff
$\alpha_1 \beta_2 = \alpha_2 \beta_1$ as polynomials in
$\mathbb{Z}[\mathcal V]$ \emph{and} $\beta_1 = \beta_2$. The
soundness of this rule rests on the following: if $\alpha_1 = c
\beta_1$ (i.e., the numerator is divisible by the denominator
exactly) then $\lfloor \alpha_1 / \beta_1 \rfloor + 1$ is exactly
$c + 1$, and the abstract value coincides with the concrete one for
every integer instantiation. When $\alpha_1$ is \emph{not}
provably divisible by $\beta_1$ on the symbolic level (the
generic case for symbolic $H, W$), \textsc{TensorGuard} attaches
a \emph{divisibility precondition}
$(\alpha_1 \mod \beta_1) = 0$ that downstream operators must
discharge before equality can be claimed. If a downstream operator
requires the floor result to be a specific concrete value (e.g.,
\verb|view(B, 16*5*5)|), the precondition becomes the
linear Diophantine constraint $\alpha_1 = 5 \beta_1$, which the
analyzer either proves (by canonical-form normalization) or fails
sound on (returning \textsf{unknown}). In the curated benchmark
all spatial inputs are concrete integers, so all floor evaluations
collapse to integer constants; in the torchvision sweep the
default input shapes ($224\times 224$, $299\times 299$,
$240\times 240$) likewise allow floor evaluation to constants and
the verifier proceeds without symbolic floor reasoning. The
fully-symbolic case is exercised only in the unit tests.

\paragraph{\texttt{view(-1, e\_2, ..., e\_m)}:} A view with one
$-1$ entry resolves the missing dimension as $\prod_i d_i / \prod_{j>1} e_j$.
Our abstract domain encodes this symbolically: we replace $-1$
with a fresh symbolic dimension $u$, emit the precondition
\[
\prod_i d_i = u \cdot \prod_{j>1} e_j,
\]
and require this product equation to be satisfied modulo
canonical-form normalization in $\mathbb Z[\mathcal V]$. If the
canonical form admits a solution in $u$ (i.e., the right-hand-side
divides the left-hand-side as polynomials), $u$ is replaced by
the quotient and propagation proceeds; otherwise the verifier
emits a divisibility-failure verdict. Concretely, on
\verb|x.view(B, -1, 32)| with $x$ of shape $(B, 16, 5, 5)$, the
constraint is $B \cdot 16 \cdot 5 \cdot 5 = u \cdot B \cdot 32$,
which canonicalises to $u = 12.5$ -- not an integer, so the
verifier emits \textsf{rejected} with witness $u = 12.5$. On
\verb|x.view(B, -1, 8)| with $x$ of shape $(B, 16, 5, 5)$, the
constraint $B \cdot 400 = u \cdot B \cdot 8$ canonicalises to
$u = 50$, an integer, and the verifier proceeds. The
divisibility check is decidable on the polynomial fragment
$\mathbb Z[\mathcal V]$ when one side is a constant times the
other; the residual cases (two non-constant polynomials) are
returned \textsf{unknown}. We do not claim completeness over
general polynomial divisibility (which would demand a
factorisation oracle for $\mathbb Z[\mathcal V]$ and is
NP-hard).

\paragraph{Soundness of the divisibility precondition.} The
divisibility test is sound by construction: if the abstract
divisibility holds, then for every concrete instantiation the
corresponding integer divisibility holds, by linearity of $\mod$
in $\mathbb Z$. Conversely, if the abstract divisibility fails,
then by the polynomial identity, there exists a concrete
instantiation that exhibits a non-zero remainder. Hence the
precondition test is sound (no false negatives) and complete on
the polynomial fragment (no false positives within the fragment).

\section{Response to reviewer concerns}\label{app:response}

We summarise the concrete revisions made in this version of the
paper in response to the reviewer's $4/10$ vote.

\textbf{(1) Retitle and reframe as fragment verifier.} The title
is now ``Static Shape Verification for a PyTorch Module Fragment
via Symbolic Abstract Interpretation'', the abstract is rewritten
to scope the tool to a DSL fragment, and the conclusion no longer
claims ``no missed errors anywhere''.

analyzed. The headline is now the union
of curated, real-repo, and injected-bug benchmarks rather than
the wrapper-only torchvision sweep.

\textbf{(3) Real-repo evaluation.} New \Cref{sec:realrepo} reports
per-file verdicts on six small public PyTorch model files
(nanoGPT, minGPT, MAE, U-Net, U-Net-parts, labml MHA), with
honest reporting of three TG false positives caused by
configuration-attribute indirection
(``\texttt{Linear out\_features unknown}'').

\textbf{(4) First-class \textsf{unknown} verdict table.} New
\Cref{tab:unknown} aggregates verdicts across all benchmarks.

\textbf{(5) Theorem language weakened, divisibility promoted.}
\Cref{thm:soundness} now states ``Sound relative to the modeled
DSL and declared opaque-module specifications'', and the
\verb|Conv2d| floor-divisibility (C1) and \verb|view|-divisibility
(C2) preconditions are explicit hypotheses of the theorem (not
buried in the appendix). The Lean development gains
\texttt{applyOp\_view\_divisible} and
\texttt{applyOp\_view\_no\_remainder}; \texttt{lake build} stays
green.

\textbf{(6) Comparison to PyTorch-native shape tools.} New
\Cref{sec:pytorch-native} (and \Cref{tab:pytorch-native-comparison})
compares against \verb|torch.fx|+ShapeProp, \verb|FakeTensorMode|,
\verb|torch.export|, and the \verb|ShapeEnv| of
\verb|torch.fx.experimental.symbolic_shapes|. We articulate
exactly what TG adds (closed-form symbolic constraints,
no-execution analysis, scoped formal soundness, friendly errors,
explicit \textsf{unknown}) and what the native ecosystem provides
that TG does not (RNN/LSTM, dynamic shapes, broad operator
coverage).

\textbf{(7) Conv2d floor / view divisibility in headline theorem.}
See (5) above; the appendix proof of \Cref{app:floor-divisibility}
is unchanged but the side conditions are now hypotheses of the
main theorem.

\textbf{(8) Injected-bug benchmark.} New \Cref{sec:injected-bug}
reports controlled-ground-truth catch rates over $10$ model
archetypes $\times$ $3$ mutations classes ($7$ mutations actually
fired); TG / FX / FakeTensor each catch $7/7$ injected bugs and
each accept $10/10$ clean baselines.

\section{Real-benchmark results}\label{app:real-benchmarks}
On the broader \texttt{real\_benchmarks/benchmark\_results.json}
harness ($44$ tests), \textsc{TensorGuard} obtains
$\mathrm{F}_1=0.7692$. The gap is dominated by out-of-fragment
operators (LSTM/GRU/custom autograd). Within the fragment the F$_1$
remains $1.0$. Extending the fragment to recurrent operators is left
to future work.

\section{Reproducibility}\label{app:repro}
A single command per experiment, all using \texttt{python3.11} and
deterministic seeds embedded in the harness scripts:
\begin{lstlisting}
cd tensorguard
# Original benchmarks
python3.11 verify_neurips.py            # base, 2.6 s, 7 models
python3.11 verify_neurips_extended.py   # ext, 5.8 s, 14 models
python3.11 verify_neurips_revision_handwritten.py
                                        # hand-written real bug
python3.11 experiments/real_repo_eval.py
                                        # 6 real public PyTorch
                                        # model files (cached)
python3.11 experiments/injected_bug_eval.py
                                        # 10 models x 3 mutations
python3.11 experiments/comparison_pytorch_native.py
                                        # TG vs FX/FakeTensor/
                                        # torch.export capability
                                        # matrix
# NEW benchmarks introduced by this revision
python3.11 benchmarks/tv_realsource_benchmark.py
                                        # NEW: 30 torchvision
                                        # (module, class) pairs;
                                        # introspecting analyzer
python3.11 benchmarks/realcode_benchmark.py
                                        # NEW: 33-file hand-authored
                                        # real-code corpus
python3.11 benchmarks/fp_ablation.py
                                        # NEW: classify FPs by cause
python3.11 benchmarks/localization_quality.py
                                        # NEW: TG vs FX vs FakeTensor
                                        # cause-line / concrete-fix
                                        # rubric on 5 buggy files
# Extended unit-test suite for the operator-catalogue extension
python3.11 -m pytest tests/test_extended_ops.py -q
                                        # 37 passing tests, 0 fail
# Lean machine-check
cd lean && lake build                   # core+extended fragment
                                        # soundness in Lean 4
\end{lstlisting}
Outputs at
\verb|experiments/neurips_validation*.json|,
\verb|experiments/neurips_revision_handwritten_bug.json|,
\verb|experiments/real_repo_eval.json|,
\verb|experiments/injected_bug_eval.json|,
\verb|experiments/comparison_pytorch_native.json|,
\verb|benchmarks/torchvision_realsource_results.json|,
\verb|benchmarks/realcode_results.json|,
\verb|benchmarks/fp_ablation_results.json|, and
\verb|benchmarks/localization_quality_results.json|. The Lean
project under \verb|lean/| pins toolchain
\texttt{leanprover/lean4:v4.14.0} via \verb|lean/lean-toolchain| and
requires no \textsc{mathlib}; \texttt{lake build} succeeds without
errors and without any \texttt{sorry}, machine-checking the
soundness theorems listed in \Cref{sec:lean-theorems}. All
deterministic; no GPU required.

\section{Lean-verified theorems}\label{sec:lean-theorems}
The Lean~4 project at \verb|lean/| machine-checks
$17$ theorems, none using \texttt{sorry}, structured as two files:
\begin{itemize}\itemsep0pt
  \item \verb|lean/TensorGuard/Soundness.lean| ($8$ theorems):
    \texttt{applyOp\_sound\_linear},
    \texttt{applyOp\_sound\_view},
    \texttt{applyOp\_view\_verdict},
    \texttt{applyOp\_view\_divisible},
    \texttt{applyOp\_view\_no\_remainder},
    \texttt{applyOp\_sound\_broadcast\_add},
    \texttt{applyOp\_linear\_rank}, and the typing-rule
    soundness lemma listed in
    \Cref{sec:soundness-proof}. These cover the
    original \texttt{Linear} / \texttt{view} /
    \texttt{broadcast\_add} core fragment.
  \item \verb|lean/TensorGuard/Extended.lean| ($9$ theorems,
    new in this revision): \texttt{matmul2\_sound},
    \texttt{bmm\_sound}, \texttt{transpose2\_involutive},
    \texttt{permList\_compose} (permute composition / double
    transpose as instance), \texttt{conv1dOut\_monotone}
    (Conv2d output spatial formula, monotonicity in the
    input spatial dim), \texttt{relu\_shape\_preserving}
    (ReLU phase preservation), \texttt{bcastDim\_comm},
    \texttt{bcast\_comm} (binary broadcast composition is
    commutative on shapes of equal rank), and
    \texttt{bcast\_rank1\_sound} (broadcast soundness on
    rank-$1$ shapes).
\end{itemize}
The extended fragment is mechanised \emph{in Lean~4 core, no
mathlib}. Each theorem is proved sorry-free in tactic mode;
\verb|lake build| terminates in $<\!10$\,s on commodity CPU. The
gap between the formally-verified fragment and the analyzer's full
operator catalogue (\Cref{app:ops-extended}) remains: see
\Cref{sec:discussion} for our scope statement.

\section{Operator coverage}\label{app:ops}
The DSL fragment $\mathcal F$ covers:
\verb|nn.Linear|, \verb|nn.Conv1d/2d/3d|, \verb|nn.BatchNorm1d/2d/3d|,
\verb|nn.LayerNorm|, \verb|nn.ReLU|, \verb|nn.GELU|, \verb|nn.Sigmoid|,
\verb|nn.Tanh|, \verb|view|, \verb|reshape|, \verb|transpose|,
\verb|permute|, \verb|cat|, \verb|stack|, \verb|squeeze|,
\verb|unsqueeze|, \verb|+|, \verb|-|, \verb|*|, \verb|/|,
\verb|torch.matmul|, \verb|@|. RNN/LSTM/GRU and custom autograd
functions are out-of-fragment (the analyzer abstains and emits an
\textsf{unknown} verdict).

\section{Extended operator catalogue}\label{app:ops-extended}
This revision wires $30$ additional operator handlers into
the introspecting analyzer
(\verb|src/tensor_shapes.py|), with $37$ unit tests in
\verb|tests/test_extended_ops.py| (all passing under
\verb|python3.11 -m pytest|). The new operators, grouped by
family:

\begin{description}\itemsep1pt
  \item[\textit{Indexing / reshape long tail.}]
    \verb|gather(dim, index)|,
    \verb|scatter|, \verb|scatter_add|, \verb|scatter_| (in-place),
    \verb|index_select(dim, index)|,
    \verb|masked_select(mask)| (returns rank-$1$, length-symbolic),
    \verb|take_along_dim|,
    \verb|narrow(dim, start, length)|,
    \verb|roll|, \verb|flip|, \verb|fliplr|, \verb|flipud|,
    \verb|rot90|,
    \verb|repeat_interleave(repeats, dim)|,
    \verb|broadcast_to(shape)|.
  \item[\textit{Convolutional family completion.}]
    \verb|nn.Conv1d|, \verb|nn.Conv3d|,
    \verb|nn.ConvTranspose2d| with the full output formula
    $H_{\textrm{out}} = (H_{\textrm{in}}-1)\cdot s -2p
    + d\cdot(k-1) + p_{\textrm{out}} + 1$,
    \verb|nn.AvgPool2d|, \verb|nn.AdaptiveMaxPool2d|.
  \item[\textit{Normalization completion.}]
    \verb|nn.BatchNorm1d|, \verb|nn.BatchNorm3d|,
    \verb|nn.GroupNorm|.
  \item[\textit{Attention.}]
    \verb|F.scaled_dot_product_attention(q, k, v)|, with
    output shape $(\ldots, L_{q}, E_{v})$, wired through
    \verb|_infer_call_shape|.
  \item[\textit{Reduction completion.}]
    \verb|var|, \verb|std|, \verb|argmax|, \verb|argmin|,
    \verb|amax|, \verb|amin|, \verb|logsumexp|, \verb|any|,
    \verb|all|, \verb|nansum|, \verb|nanmean|,
    \verb|count_nonzero|, all honouring the \verb|dim| and
    \verb|keepdim| arguments.
\end{description}

The soundness statement for the new handlers is the
\emph{abstain-on-unknown} rule that the existing
\textsc{TensorGuard} pipeline already enforces: a transfer rule
returns \texttt{None} (rather than fabricate a shape) whenever any
of its arguments is symbolic in a way the rule cannot resolve.
This is verified by the \verb|test_extended_ops.py| suite, which
includes $5$ bug-injection tests
(\verb|Linear|/\verb|Conv1d|/\verb|Conv2d| input-channel
mismatch, \verb|ConvTranspose2d| downstream \verb|Linear|
mismatch, \verb|matmul| inner-dim mismatch) that the analyzer
correctly catches, alongside $32$ correctness tests on
in-fragment usage that report no false positives.

The operator catalogue currently does \emph{not} introduce new
shape-violation \emph{checks} for the indexing family
(\verb|gather|, \verb|scatter|, \dots); the new handlers infer
output shape only, deferring violation detection to the existing
\verb|model_checker| machinery for the operators it already
supports. This is conservative-by-design and is honest about the
scope of the catalogue extension.

\section{Wrapper-only torchvision sweep (archive)}
\label{app:tv-wrapper}
For backward compatibility with the previous version of this
paper, we report the original wrapper-only torchvision sweep here.
$30$ \verb|nn.Module| wrappers, each of the form
\verb|self.net = tvm.<arch>(weights=None); return self.net(x)|, are
verified by treating \verb|self.net(x)| as an opaque submodule call
with the standard $\textrm{rank}(\mathrm{in})=4$,
$\textrm{rank}(\mathrm{out})=2$ specification.
\textsc{TensorGuard} returns $30/30$ \textsf{verified}, mean
$0.42$\,s/model; \verb|torch.fx| ShapeProp returns $30/30$ ok at
roughly $3.5$\,s/model; \verb|FakeTensorMode| forward succeeds on
$28/30$ (the two failures are inside FakeTensor's RegNet dispatch).
This sweep is, we now agree, an \emph{opaque-boundary sanity
check}: it shows TG's modular abstention does not blow up at the
library boundary, but it does not demonstrate any actual
introspection inside the library subtree. The real-source
introspection benchmark of \Cref{sec:tv-source} is the version of
this experiment that probes the analyzer's actual capability.

\end{document}
